\documentclass[11pt]{amsart}

\usepackage{amsmath,amssymb,mathtools}
\usepackage{booktabs}
\usepackage{microtype}
\usepackage{xcolor}
\usepackage{xurl}
\usepackage[colorlinks=true,linkcolor=blue!45!black,citecolor=blue!45!black,
  urlcolor=blue!45!black]{hyperref}
\urlstyle{tt}
\emergencystretch=3em

\newcommand{\K}{\mathbb{K}}
\newcommand{\Q}{\mathbb{Q}}
\newcommand{\Sym}{\operatorname{Sym}}
\newcommand{\Span}{\operatorname{span}}
\newcommand{\im}{\operatorname{im}}
\newcommand{\rank}{\operatorname{rank}}
\newcommand{\essdim}{\operatorname{essdim}}
\newcommand{\perm}{\operatorname{perm}}
\newcommand{\CR}{\mathrm{R}_{\mathrm{Chow},\K}}
\newcommand{\cD}{\mathcal{D}}

\newtheorem{theorem}{Theorem}[section]
\newtheorem{proposition}[theorem]{Proposition}
\newtheorem{lemma}[theorem]{Lemma}
\theoremstyle{definition}
\newtheorem{definition}[theorem]{Definition}
\newtheorem{remark}[theorem]{Remark}

\title[Certified Local Computations for the Six-by-Six Permanent]
{Certified Local Computations and the Ordinary Chow Rank of the
Six-by-Six Permanent}
\author{Rongyu Xu}
\hypersetup{
  pdftitle={Certified Local Computations and the Ordinary Chow Rank of the Six-by-Six Permanent},
  pdfauthor={Rongyu Xu}}
\date{September 1, 2026}
\subjclass[2020]{Primary 15A69; Secondary 14N05, 68V05}
\keywords{permanent, Chow rank, product rank, catalecticant,
computer-assisted proof}

\begin{document}

\begin{abstract}
Let \(\perm_6\) be the permanent of a generic \(6\) by \(6\) matrix.
Over an algebraically closed field of characteristic zero, we prove that its
ordinary Chow rank is \(32\). This submission edition isolates the finite
local computations as explicit lemmas with declared inputs, exhaustive
domains, exact outputs, and verifier boundaries. The lower bound combines a
symmetric image-span inequality with a half-defect estimate for arbitrary
quotient derivative symbols of a sextic Chow term. A primary exact replay
reconstructs all finite tables. A second implementation independently checks
every labelled bipartite graph with at most six edges and all \(45{,}696\)
coordinate squarefree-symbol cases. No floating-point or randomized step is
used. A separate standard-library implementation independently expands and
differentiates all five dependent-factor normal forms. No border-rank or
general-\(n\) equality is claimed.
\end{abstract}

\maketitle

\section{Introduction}

The Chow variety in degree six parametrizes products of six linear forms.
Its secant geometry leads to two distinct invariants: ordinary Chow rank,
which requires an actual finite sum of products, and border Chow rank, which
allows limits of such sums.  This article concerns only the ordinary invariant.
Product rank is synonymous with ordinary Chow rank in the relevant literature;
see \cite{IltenTeitler2016}.

Glynn's identity gives a \(32\)-term product decomposition of \(\perm_6\).
The content of Theorem~\ref{thm:main} is the matching lower bound.  The proof
uses the middle-to-quadratic derivative tower, a small-intersection theorem
for permanent quadrics, and a local half-defect estimate for every sextic Chow
term.  A symmetric image-span inequality and a filtration of the factor spans
give two estimates with the same middle-rank defect term. Section~\ref{sec:global}
defines every quantity in those estimates; cancellation of their common
defect forces at least \(32\) summands.

The local half-defect estimate contains the finite part of the argument.  It
uses exact rational differentiation of the five dependent-factor normal
forms, an exact enumeration of \(45{,}696\) coordinate symbol cases, and a
coordinate audit of the small permanent-quadratic intersections.  The replay
reconstructs these objects from their definitions and compares them with the
frozen payload.  All geometric reductions and global inequalities are written
proofs.  The final status is an author-reviewed computer-assisted theorem;
named external peer review and proof-assistant formalization remain outside
the claim.

\section{Statement and permanent derivatives}

\subsection*{Notation and conventions}

Let \(\K\) be an algebraically closed field of characteristic zero. For a
positive integer \(n\), write \([n]=\{1,2,\ldots,n\}\) and let \(S_n\) be the
group of bijections of \([n]\). If \(W\) is a finite-dimensional
\(\K\)-vector space, then \(W^*\), \(\Sym^jW\), and \(\operatorname{Gr}(a,W)\)
denote its dual, its degree-\(j\) symmetric power, and the Grassmannian of
\(a\)-dimensional subspaces. For a finite set \(\Omega\),
\(\binom{\Omega}{a}\) denotes the set of its \(a\)-element subsets. For a
subspace \(H\subset W\), its annihilator is
\(H^\perp:=\{\lambda\in W^*: \lambda(h)=0\text{ for every }h\in H\}\).
We use \(\Span\), \(\im\), \(\ker\), \(\rank\), and \(\dim\) with their
usual linear-algebraic meanings.

Let
\[
 V=\Span_{\K}\{x_{rc}:1\le r,c\le6\},\qquad \dim V=36.
\]
Write \(X=(x_{rc})_{r,c\in[6]}\) and define
\[
 \perm_6=\sum_{\sigma\in S_6}\prod_{r=1}^6x_{r,\sigma(r)}\in\Sym^6V.
\]
For \(f\in\Sym^mV\), \(0\le a\le m\), and
\(\lambda_1,\ldots,\lambda_a\in V^*\), let
\(\partial_{\lambda_1}\cdots\partial_{\lambda_a}f\) denote iterated
directional differentiation. With \(j=m-a\), define
\[
 \begin{aligned}
 C_{a,j}(f)&:\Sym^aV^*\longrightarrow\Sym^jV,\\
 C_{a,j}(f)(\lambda_1\cdots\lambda_a)
   &:=\partial_{\lambda_1}\cdots\partial_{\lambda_a}f,\\
 \cD_j(f)&:=\im C_{a,j}(f).
 \end{aligned}
\]
For \(G\subset\Sym^jV\), put
\(\partial G:=\Span\{\partial_\lambda g:\lambda\in V^*,\ g\in G\}\), and
write \(\partial g:=\partial\Span\{g\}\). For \(g\in\Sym^jV\), let \(\operatorname{Ess}(g)\) be the smallest subspace
\(L\subset V\) for which \(g\in\Sym^jL\). Its dimension is the number of
essential variables of \(g\). For \(q\in\Sym^2V\), let \(M(q)\) be the
symmetric matrix of its associated bilinear form in the ordered basis
\((x_{rc})\), and define the quadratic rank of \(q\) to be \(\rank M(q)\).
Finally, put
\(E_j:=\cD_j(\perm_6)\).

For \(I,J\subset[6]\) with \(|I|=|J|\), let \(X_{I,J}\) denote the
corresponding submatrix of \(X\), and let \(\perm(X_{I,J})\) denote its
permanent, defined by the same permutation formula as above.

\begin{definition}
A sextic Chow term is a nonzero product \(T=\ell_1\cdots\ell_6\). The
ordinary Chow rank \(\CR(f)\) of \(f\in\Sym^6V\) is the least positive
integer \(N\) for which \(f=T_1+\cdots+T_N\) with sextic Chow terms
\(T_1,\ldots,T_N\); we set \(\CR(0)=0\).
\end{definition}

\begin{theorem}\label{thm:main}
Over \(\K\),
\[
 \boxed{\CR(\perm_6)=32.}
\]
\end{theorem}

We first record the permanent derivative facts used below.

\begin{lemma}[Permanent derivative tower]\label{lem:tower}
One has
\[
 \dim E_3=400,\qquad \dim E_2=225,\qquad
 \partial E_2=V,
\]
and
\[
 E_2^{(1)}:=\{g\in\Sym^3V:\partial g\subset E_2\}=E_3. \tag{1.1}
\]
Moreover, every nonzero element of \(E_2\) has quadratic matrix rank at
least four, and every nonzero element of \(E_3\) has at least nine essential
variables.
\end{lemma}

\begin{proof}
The \(3\times3\) and \(2\times2\) subpermanents are bases of \(E_3\) and
\(E_2\): distinct row-column blocks have disjoint matching supports. Their
numbers are \(\binom63^2=400\) and \(\binom62^2=225\). Differentiating a
rectangle permanent at one corner gives the opposite variable, proving
\(\partial E_2=V\).

The inclusion \(E_3\subset E_2^{(1)}\) is immediate. Conversely, suppose
all first derivatives of a cubic \(g\) lie in \(E_2\). Row and column
multiaffinity forces every monomial of \(g\) to be a three-edge matching.
Inside a fixed \(3\times3\) block, the rectangle relations equate
coefficients of matching monomials related by a transposition. Transpositions
connect all six matchings, so each block is a scalar subpermanent. This proves
(1.1).

The row-column diagonal torus has distinct characters on each subpermanent
basis. A generic one-parameter subgroup degenerates any nonzero element to a
single basis element, while matrix and first-catalectic ranks cannot increase
under specialization. A rectangle permanent has quadratic rank four. The
nine first derivatives of a \(3\times3\) permanent are independent. The two
rank floors follow.
\end{proof}

\section{Small permanent-quadratic intersections}

\begin{lemma}\label{lem:small-intersection}
If \(L\subset V\) has dimension at most six, then
\[
 \dim(E_2\cap\Sym^2L)\le3.
\]
If \(\dim L\le5\), the upper bound improves to one.
\end{lemma}

\begin{proof}
For fixed nonnegative integers \(a,q\), the locus in
\(\operatorname{Gr}(a,V)\) on which
\(\dim(E_2\cap\Sym^2L)\ge q\) is closed, projective, and stable under the
connected row-column torus. If nonempty, the Borel fixed-point theorem gives
a torus-fixed point \cite[Chapter~I, Theorem~4.7]{Borel1991}. Since the
variable weight spaces are one-dimensional,
such a point is a coordinate \(a\)-plane, equivalently a bipartite graph with
\(a\) edges.

At a coordinate point, the intersection has one independent rectangle
permanent for each four-cycle. Two distinct four-cycles have a union of at
least six edges. If the union has exactly six edges, it is \(K_{2,3}\) or
\(K_{3,2}\), which has exactly three four-cycles. Thus five edges support at
most one four-cycle and six edges support at most three.
\end{proof}

\section{Local linear algebra}\label{sec:local-linear-algebra}

If \(L\subset V\) and \((y_1,\ldots,y_\ell)\) is a basis with dual basis
\((y_1^*,\ldots,y_\ell^*)\), define the polarization map
\[
 \delta_L:\Sym^3L\longrightarrow L\otimes\Sym^2L,
 \qquad \delta_L(g)=\sum_{j=1}^{\ell}y_j\otimes\partial_{y_j^*}g.
\]
This definition is basis-independent and \(\delta_L\) is injective. We write
\(\delta\) when \(L\) is clear.

\begin{lemma}[Kernel-preimage bound]\label{lem:kernel-preimage}
Let \(D\) be a vector space, let \(P:L\twoheadrightarrow D\) be a surjective
linear map of rank \(d=\dim D\), let \(S:=\ker P\), and let
\(R\subset\Sym^2L\). If \(1_{\Sym^2L}\) denotes the identity of
\(\Sym^2L\), then
\[
 \ker((P\otimes1_{\Sym^2L})\delta)=\Sym^3S,
\]
and the inverse image of \(D\otimes R\) under
\((P\otimes1_{\Sym^2L})\delta\) has dimension at most
\[
 \boxed{\binom{\dim L-d+2}{3}+d\dim R.} \tag{3.1}
\]
\end{lemma}

\begin{proof}
The first equality follows in coordinates: all derivatives in directions
modulo \(S\) vanish precisely when the cubic depends only on \(S\). The
inverse image of a subspace has dimension at most the kernel dimension plus
the subspace dimension, which gives (3.1).
\end{proof}

\begin{lemma}[Squarefree projected-symbol table]\label{lem:squarefree}
Let \(\widehat L=\K^6\) with basis \(z_1,\ldots,z_6\), and define the
squarefree quadratic and cubic spaces by
\[
 \begin{aligned}
 \widehat F&:=\Span\{z_az_b:1\le a<b\le6\},\\
 \widehat U&:=\Span\{z_az_bz_c:1\le a<b<c\le6\}.
 \end{aligned}
\]
The restriction of polarization to \(\widehat U\) is determined by
\[
 \delta(z_az_bz_c)=z_a\otimes z_bz_c+z_b\otimes z_az_c
                    +z_c\otimes z_az_b.
\]
If \(P:\widehat L\twoheadrightarrow D\) has rank \(d\) and
\(R\subset\widehat F\) has dimension at most \(r\le3\), let
\(q_R:\widehat F\twoheadrightarrow\widehat F/R\) be the quotient map. Then the following
table is a rowwise lower bound for
\(\rank((P\otimes q_R)\delta)\):
\[
\begin{array}{c|rrrrrrr}
r\backslash d&0&1&2&3&4&5&6\\ \hline
0&0&10&16&19&20&20&20\\
1&0&9&14&16&16&20&20\\
2&0&8&12&13&16&19&20\\
3&0&7&10&10&15&17&19.
\end{array} \tag{3.2}
\]
\end{lemma}

\begin{proof}
The diagonal torus preserves \(\delta\). The orbit closure of
\((\ker P,R)\) in the relevant product of Grassmannians contains a fixed
point, and rank can only drop in specialization. Thus \(\ker P\) may be
taken as the coordinate span indexed by a vertex set \(B\subset[6]\), and
\(R\) as the span of monomials indexed by an edge set
\(A\subset\binom{[6]}2\) with at most \(r\) edges. Put
\(O:=[6]\setminus B\). Let \(e_A(B)\) be the number of edges of \(A\) with
both endpoints in \(B\), let \(\deg_A(b,O)\) be the number of neighbors of
\(b\) in \(O\), and let \(A[O]\) be the graph induced by \(A\) on \(O\).

Distinct source cubics have disjoint output supports. The number of killed
cubics is
\[
 \binom{|B|}{3}+|O|e_A(B)
 +\sum_{b\in B}\binom{\deg_A(b,O)}2
 +\#\{\text{triangles of }A[O]\}. \tag{3.3}
\]
Inspection of edge sets of size at most three gives (3.2). The exact replay
independently checks all 45,696 coordinate pairs and records a minimizer for
each entry.
\end{proof}

\section{The half-defect estimate}\label{sec:half-defect}

Let \(T=\ell_1\cdots\ell_6\ne0\), and set
\[
 L=\Span\{\ell_1,\ldots,\ell_6\},\quad
 U=\cD_3(T),\quad u=\dim U,
\]
\[
 F=\cD_2(T),\quad R=F\cap E_2.
\]
For a surjection \(P:L\twoheadrightarrow D\) onto a vector space \(D\) of
dimension \(d\), let
\(q_R:F\twoheadrightarrow F/R\) be the quotient map and define
\[
 \beta_{P,R}:U\xrightarrow{\delta}L\otimes F
 \xrightarrow{P\otimes q_R}D\otimes(F/R).
\]

\begin{proposition}[Half-defect quotient-symbol estimate]
\label{prop:half-defect}
For every nonzero sextic Chow term and every quotient of its actual factor
span,
\[
 \boxed{\rank\beta_{P,R}+\frac{20-u}{2}\ge\frac{10}{3}d.} \tag{4.1}
\]
\end{proposition}

\begin{proof}
Write \(\ell=\dim L\).
We call \(\beta_{P,R}\) a \emph{full symbol} when \(d=\ell\), equivalently
when \(P\) is an isomorphism. Whenever a parenthesized row is displayed below,
its entries are indexed from left to right by \(d=0,1,\ldots,\ell\).

\paragraph{Factor span at most four.}
Choose \(\ell\) independent factors as coordinates. A generic diagonal
one-parameter subgroup degenerates the remaining factors to unique coordinate
initials, so \(T\) specializes to
\[
 x_1^{m_1}\cdots x_\ell^{m_\ell},\qquad m_j\ge1,
 \qquad\sum_jm_j=6.
\]
Middle rank cannot increase in the limit. Counting degree-three divisors over
the positive partitions of six gives
\[
 \ell=1,2,3,4\quad\Longrightarrow\quad u\ge1,2,4,8. \tag{4.2}
\]
Lemma~\ref{lem:small-intersection} gives \(\dim R\le1\). For
\(0<d<\ell\), Lemma~\ref{lem:kernel-preimage} gives
\[
 \rank\beta_{P,R}\ge
 \max\left\{0,u-\binom{\ell-d+2}{3}-d\right\}. \tag{4.3}
\]
At \(d=\ell\), a kernel cubic lies in \(E_2^{(1)}=E_3\), but it is
supported on at most six essential variables, contrary to
Lemma~\ref{lem:tower}. Thus the full symbol is injective. Minimizing over the
conservative interval from (4.2) to
\(\min\{20,\binom{\ell+2}{3}\}\) gives
\[
\begin{array}{c|c}
\ell&\rank\beta_{P,R}+(20-u)/2\text{ for }d=0,\ldots,\ell\\ \hline
4&(0,9/2,8,10,14)\\
3&(5,15/2,9,12)\\
2&(8,9,11)\\
1&(19/2,21/2).
\end{array} \tag{4.4}
\]

\paragraph{Six independent factors.}
Here \(u=20\), the actual derivative spaces are the squarefree spaces, and
\(\dim R\le3\). The last row of (3.2), improved at \(d=6\) by full-symbol
injectivity, gives
\[
 (0,7,10,10,15,17,20). \tag{4.5}
\]

\paragraph{Five-dimensional factor span.}
The unique relation among the six factors puts \(T\), after changing and
rescaling coordinates, into exactly one form
\[
 T_s=x_1x_2x_3x_4x_5(x_1+\cdots+x_s),\qquad1\le s\le5. \tag{4.6}
\]
Lemma~\ref{lem:small-intersection} gives \(\dim R\le1\). Exact
differentiation of the actual polynomial gives
\[
\begin{array}{c|rrrrr}
s&1&2&3&4&5\\ \hline
\dim F&11&11&13&14&15\\
u&14&14&18&20&20.
\end{array} \tag{4.7}
\]
The replay constructs the order-four and order-three derivative matrices of
(4.6) in the ordinary monomial bases and row-reduces them over \(\mathbb Q\).

For \(s=3\), the actual space \(U\) contains the nine squarefree cubics
other than \(x_1x_2x_3\). For \(s=4,5\), it contains all ten squarefree
cubics. These subspaces are stable under the diagonal torus. A nonzero
rank-one quotient specializes to a coordinate functional, whose rank is the
number of displayed triples through that coordinate: at least five for
\(s=3\) and six for \(s=4,5\). Quotienting by \(R\) loses at most one.
Together with (3.1) and full-symbol injectivity this gives
\[
\begin{array}{c|c|c}
s&u&\rank\beta_{P,R}+(20-u)/2\text{ for }d=0,\ldots,5\\ \hline
4,5&20&(0,5,8,13,15,20)\\
3&18&(1,5,7,12,14,19).
\end{array} \tag{4.8}
\]

For \(s=1\), the actual \(U\) is spanned by the ten squarefree cubics and
\(x_1^2x_j\), \(2\le j\le5\). If a direction
\(\lambda=\sum a_j\partial_{x_j}\) has \(a_j\ne0\) for some \(j\ge2\),
the six squarefree cubics divisible by \(x_j\) and \(x_1^2x_j\), projected
modulo \(x_jL\), give seven independent quadrics. The pure \(x_1\)
direction has rank ten.

For \(s=2\), a basis consists of the seven squarefree cubics containing at
most one of \(x_1,x_2\), the six cubics
\[
 x_2^2x_j+2x_1x_2x_j,\qquad
 x_1x_2x_j+\tfrac12x_1^2x_j\quad(3\le j\le5),
\]
and \(x_1x_2(x_1+x_2)\). If \(a_j\ne0\) for some \(j\ge3\), five
squarefree inputs and the two special inputs give, modulo \(x_jL\), five
distinct pair monomials and the independent quadrics
\(x_2^2+2x_1x_2\) and \(x_1x_2+\tfrac12x_1^2\). If
\(\lambda=a_1\partial_{x_1}+a_2\partial_{x_2}\), the image contains the
three pair monomials on \(x_3,x_4,x_5\), a nonzero direction in each
\(\Span\{x_1x_j,x_2x_j\}\), and the nonzero quadratic
\[
 a_2x_1^2+2(a_1+a_2)x_1x_2+a_1x_2^2.
\]
Thus every nonzero directional map has rank at least seven, and every
positive-rank quotient symbol has rank at least six after reducing by \(R\).

For \(d=4\), let \(S=\ker P\) be a line and let \(v\) lie in the symbol
kernel. Its four complementary derivatives span a subspace of \(R\). If
that span were nonzero, the essential-variable space of \(v\) would have
dimension at most two and would support a nonzero element of \(E_2\),
contrary to the rank-four floor in Lemma~\ref{lem:tower}. Hence
\(v\in\Sym^3S\), so the kernel has dimension at most one. The full symbol at
\(d=5\) is injective. With (3.1), the \(s=1,2\) row is
\[
 (3,9,9,10,16,17). \tag{4.9}
\]

Every entry of (4.4), (4.5), (4.8), and (4.9) is at least \(10d/3\).
This proves (4.1).
\end{proof}

\begin{remark}[The repaired boundary]\label{rem:repair}
The formal pair- and triple-product spans for \(T_2\) have dimensions 12
and 17, while the actual derivative spaces have dimensions 11 and 14. For
\(T_3\), the formal triple-product span has dimension 19 while the actual
middle derivative space has dimension 18. The proof above deliberately uses
the formal squarefree table only in the six-independent-factor case. The
negative equalities are frozen as regression tests.
\end{remark}

\section{A symmetric image-span inequality}

\begin{lemma}\label{lem:image-span}
Let \(A_1,\ldots,A_N:W^*\to W\) be linear maps that are symmetric under the
natural pairing between \(W^*\) and \(W\). Set \(r_i:=\rank A_i\) and
\(D:=\dim(\sum_{i=1}^N\im A_i)\). Then
\[
 \rank\Bigl(\sum_iA_i\Bigr)\ge2D-\sum_i r_i. \tag{5.1}
\]
\end{lemma}

\begin{proof}
Choose an injective map \(B_i:\K^{r_i}\to W\) with image \(\im A_i\), and
let \(B_i^*:W^*\to(\K^{r_i})^*\) be its dual map. Symmetry gives
\(\ker A_i=(\im A_i)^\perp\), so
\(A_i=B_iJ_iB_i^*\) for an invertible symmetric map
\(J_i:(\K^{r_i})^*\to\K^{r_i}\). Let
\(B:\bigoplus_i\K^{r_i}\to W\) be the map with blocks \(B_i\), and let
\(J:=\bigoplus_iJ_i\). Sylvester's inequality
gives
\[
 \rank(BJB^*)\ge\rank(BJ)+\rank(B^*)-\sum_i r_i
 =2D-\sum_i r_i.
\]
\end{proof}

\section{The global lower bound}\label{sec:global}

Suppose that \(N\) is a positive integer and
\[
 \perm_6=\sum_{i=1}^NT_i
\]
with no zero summands. Let
\[
 A_i=C_{3,3}(T_i),\quad U_i=\im A_i,\quad u_i=\dim U_i,
 \quad\delta_i=20-u_i,\quad\Delta=\sum_i\delta_i.
\]
Put \(U=\sum_iU_i\) and \(\dim U=400+h\). The sum of the symmetric maps
\(A_i\) is the permanent middle catalecticant, with rank 400 and image
\(E_3\). Lemma~\ref{lem:image-span} gives
\[
 \boxed{h\le10N-200-\frac\Delta2.} \tag{6.1}
\]

Let \(F_i:=\cD_2(T_i)\), \(F:=\sum_iF_i\), and \(Q:=F/E_2\), and let
\(q:F\twoheadrightarrow Q\) be the quotient map. Linearity of the
derivative spaces gives \(E_2\subset F\). If
\(L_i:=\Span\{\ell_{i1},\ldots,\ell_{i6}\}\) is the actual factor span of
\(T_i=\ell_{i1}\cdots\ell_{i6}\), define
\[
 \widetilde\beta:\bigoplus_{i=1}^NU_i\longrightarrow V\otimes Q,
 \qquad (g_i)_{i=1}^N\longmapsto
 \sum_{i=1}^N(\iota_i\otimes q|_{F_i})\delta_{L_i}(g_i),
\]
where \(\iota_i:L_i\hookrightarrow V\) is inclusion.
Its kernel consists precisely of tuples for which every first derivative of
\(\sum_i g_i\) lies in \(E_2\), equivalently
\(\sum_i g_i\in E_3\) by Lemma~\ref{lem:tower}. Since \(E_3\subset U\),
\[
 \rank\widetilde\beta=\dim(U/E_3)=h. \tag{6.2}
\]

Since
\(E_2\subset F\), \(\partial E_2=V\), and \(\partial F_i\subset L_i\),
the spaces \(L_i\) span all 36 variables. Choose any ordering and put
\[
 W_0:=0,\qquad W_i:=\sum_{j=1}^iL_j,
 \qquad d_i=\dim(W_i/W_{i-1}),
 \qquad\sum_i d_i=36.
\]

Let
\(Z_i:=\im((1_{L_i}\otimes q|_{F_i})\delta_{L_i})\subset L_i\otimes Q\),
where \(1_{L_i}\) is the identity map of \(L_i\). Set
\(H_0:=0\) and \(H_i:=\sum_{j=1}^iZ_j\). The quotient projection
\[
 \pi_i:W_i\otimes Q\longrightarrow(W_i/W_{i-1})\otimes Q
\]
kills \(H_{i-1}\), hence
\[
 \dim H_i-\dim H_{i-1}\ge\dim\pi_i(H_i)=\dim\pi_i(Z_i). \tag{6.3}
\]
On the current block, the first-factor map is
\[
 P_i:L_i\twoheadrightarrow L_i/(L_i\cap W_{i-1}),
 \qquad\rank P_i=d_i.
\]
The natural map
\[
 F_i/(F_i\cap E_2)\longrightarrow F/E_2
\]
is injective. Therefore \(\dim\pi_i(Z_i)\) is exactly the local symbol rank
in Proposition~\ref{prop:half-defect}. Summing (6.3) gives
\[
 \boxed{h\ge\sum_i\left(\frac{10}{3}d_i-\frac{\delta_i}{2}\right)
 =120-\frac\Delta2.} \tag{6.4}
\]
Comparing (6.1) and (6.4) cancels the complete individual middle-rank defect
and yields \(N\ge32\).

For the reverse inequality, Glynn's identity \cite{Glynn2010} gives
\[
 \perm_6=\frac1{32}
 \sum_{\epsilon_1=1,\ \epsilon_2,\ldots,\epsilon_6=\pm1}
 \left(\prod_{r=1}^6\epsilon_r\right)
 \prod_{c=1}^6\left(\sum_{r=1}^6\epsilon_rx_{rc}\right).
\]
Walsh cancellation leaves exactly the monomials that use each row once. This
is a 32-term Chow decomposition. Theorem~\ref{thm:main} follows.

\section{Computational lemmas and certificate contract}
\label{sec:certificates}

A \emph{certificate contract} in this section consists of a finite input
class, a specified exact arithmetic domain, an asserted output, and a verifier
that exits unsuccessfully unless it reconstructs that output. Vertices of the
complete graph \(K_6\) and the two color classes of the bipartite graph
\(K_{6,6}\) are both ordered by \([6]\). Subsets and monomials are ordered
lexicographically. All ranks are exact ranks over \(\Q\).

\begin{lemma}[Coordinate intersection certificate]
\label{lem:certificate-intersection}
Among coordinate subspaces spanned by \(e\) variables of \(V\), for
\(0\le e\le6\), the maximum number of \(2\)-by-\(2\) permanent quadrics
supported on the subspace is
\[
 (0,0,0,0,1,1,3).
\]
\end{lemma}

\begin{proof}[Certificate proof]
A coordinate subspace is identified with an \(e\)-edge subgraph of
\(K_{6,6}\), and a supported rectangle permanent is identified with a
four-cycle. The primary replay enumerates row-permutation representatives.
The independent program \path{scripts/n6_independent_finite_core_audit.py}
instead enumerates every labelled edge subset, namely
\(\sum_{e=0}^6\binom{36}{e}=2{,}391{,}496\) graphs, and directly counts pairs
of common neighbors. Both computations give the displayed maxima.
\end{proof}

\begin{lemma}[Squarefree-symbol certificate]
\label{lem:certificate-squarefree}
Let \(0\le r\le3\) and \(0\le d\le6\) have the meanings fixed in
Section~\ref{sec:local-linear-algebra}: \(r\) bounds the number of killed edge
coordinates and \(d\) is the quotient dimension. The minimum coordinate
squarefree-symbol ranks are
\[
\begin{array}{c|rrrrrrr}
r\backslash d&0&1&2&3&4&5&6\\ \hline
0&0&10&16&19&20&20&20\\
1&0&9&14&16&16&20&20\\
2&0&8&12&13&16&19&20\\
3&0&7&10&10&15&17&19
\end{array}
\]
and the exhaustive coordinate domain contains \(45{,}696\) cases.
\end{lemma}

\begin{proof}[Certificate proof]
For each coordinate kernel and each killed-edge set of cardinality at most
\(r\), the verifier visits all twenty three-subsets of \([6]\). A cubic
coordinate survives exactly when one of its three vertex-edge factorizations
survives. The primary replay also evaluates the closed killed-column formula;
the independent audit uses only this direct survival test and does not import
the primary module or read the frozen JSON. It reproduces the table and the
case count exactly.
\end{proof}

\begin{lemma}[Normal-form and half-defect certificate]
\label{lem:certificate-half-defect}
For the five dependent-factor normal forms in
Section~\ref{sec:half-defect}, the exact derivative-dimension pairs are
\[
 (11,14),(11,14),(13,18),(14,20),(15,20).
\]
The exact low-factor-span floors are \(1,2,4,8\), and all half-defect rows
listed in Section~\ref{sec:half-defect} satisfy the required inequality.
Consequently the finite
local input to the cancellation in Section~\ref{sec:global} is valid.
\end{lemma}

\begin{proof}[Certificate proof]
The primary verifier expands each normal form as a sparse polynomial with
rational coefficients, differentiates it symbolically, and row-reduces the
actual derivative coefficient matrices. It separately enumerates the positive
partitions of six into at most four parts to obtain the monomial floors, then
evaluates every displayed half-defect row as a rational number. In particular,
it tests actual derivative spaces rather than substituting formal pair- or
triple-product spans. A missing profile, an unequal dimension, or a failed
inequality terminates the replay with an error. This lemma is independently
specified here and is also checked by the standard-library program
\path{scripts/n6_dependent_normal_forms_independent.py}, which does not import
the primary verifier or read its frozen JSON.
\end{proof}

These three lemmas state the mathematical content supplied by computation.
The fixed-point reductions, prolongation identity, directional cases,
symmetric image-span inequality, filtration increment, and defect cancellation
remain written arguments in the preceding sections.

\section{Exact replay and audit boundary}

From the repository checkout, run
\begin{verbatim}
python scripts/n6_exact_ordinary_chow_rank_32.py \
  --verify-json data/n6_exact_ordinary_chow_rank_32.json
python -m unittest \
  tests.test_n6_exact_ordinary_chow_rank_32 -v
python scripts/n6_independent_finite_core_audit.py
python scripts/n6_dependent_normal_forms_independent.py \
  --output evidence/n6_dependent_normal_forms_independent.json
\end{verbatim}
All four commands require Python 3 and only its standard library. The first
command derives the coordinate intersection maxima, all 45,696
coordinate cases behind (3.2), the actual normal-form derivative ranks and
squarefree subspaces in (4.7)--(4.8), the low-span monomial floors, and the
final cancellation. It also records the formal/actual mismatches in
Remark~\ref{rem:repair}.

The unit tests exercise the primary implementation as regression tests. The
third command is independent of that implementation and of the frozen JSON;
its scope is exactly Lemmas~\ref{lem:certificate-intersection}
and~\ref{lem:certificate-squarefree}. The fourth command independently
reconstructs the normal-form derivative ranks used in
Lemma~\ref{lem:certificate-half-defect}.

The public proof package places the exact data, primary verifier, independent
verifiers, tests, and generated receipts beside this manuscript so that no
private repository access is required.

The Borel fixed-point reductions, prolongation identity, directional
arguments, symmetric image-span lemma, and filtration increment (6.3) are
written proofs rather than numerical inferences. The result is an exact
computer-assisted theorem reviewed by the author. Named external peer review
and proof-assistant formalization remain outstanding.

\begin{remark}[Scope]
The theorem concerns ordinary Chow rank over characteristic zero. It does not
determine border Chow rank or prove the analogous equality
\(\mathrm{R}_{\mathrm{Chow},\K}(\perm_n)=2^{n-1}\) for the permanent
\(\perm_n\) of a generic \(n\) by \(n\) matrix. Product rank and Chow rank are
used synonymously in the relevant
literature; see \cite{IltenTeitler2016}.
\end{remark}

\begin{thebibliography}{9}

\bibitem{Borel1991}
A.~Borel,
\emph{Linear algebraic groups}, second ed., Graduate Texts in Mathematics,
vol.~126, Springer-Verlag, New York, 1991,
\href{https://doi.org/10.1007/978-1-4612-0941-6}
{doi:10.1007/978-1-4612-0941-6}.

\bibitem{Glynn2010}
D.~G. Glynn,
\emph{The permanent of a square matrix},
European J. Combin. \textbf{31} (2010), no.~7, 1887--1891,
\href{https://doi.org/10.1016/j.ejc.2010.01.010}
{doi:10.1016/j.ejc.2010.01.010}.

\bibitem{IltenTeitler2016}
N.~Ilten and Z.~Teitler,
\emph{Product ranks of the \(3\times3\) determinant and permanent},
Canad. Math. Bull. \textbf{59} (2016), no.~2, 311--319,
\href{https://doi.org/10.4153/CMB-2015-076-1}
{doi:10.4153/CMB-2015-076-1}.

\end{thebibliography}

\end{document}
